% !TEX program = xelatex
% c5p-s-2026-final-answer.tex
% 2026 €€£ 超专业级别硬件能力认证第一轮（C5P-S1）提高级 C++语言试题·参考答案与解析
\documentclass[a4paper]{ctexart}

\usepackage{fontspec}
\usepackage{geometry}
\geometry{top=2.7cm, left=1.9cm, right=1.9cm, bottom=2.9cm}
\setmainfont{Consolas}
\setCJKmainfont[BoldFont=SimHei]{SimSun}
\usepackage{tikz}
\usepackage{eso-pic}
\newcommand{\waterword}{%
  \makebox[0pt][c]{%
    \tikz[baseline]{%
      \node[inner sep=0pt, rotate=-35, text=black, opacity=0.085,
        font=\fontsize{54}{70}\selectfont] {折工};%
    }%
  }%
}
\AddToShipoutPictureFG{%
  \put(\LenToUnit{-0.3cm},\LenToUnit{1.2cm}){\waterword}%
  \put(\LenToUnit{5.0cm},\LenToUnit{0.0cm}){\waterword}%
  \put(\LenToUnit{10.3cm},\LenToUnit{1.4cm}){\waterword}%
  \put(\LenToUnit{15.6cm},\LenToUnit{0.2cm}){\waterword}%
  \put(\LenToUnit{2.2cm},\LenToUnit{6.2cm}){\waterword}%
  \put(\LenToUnit{7.5cm},\LenToUnit{4.8cm}){\waterword}%
  \put(\LenToUnit{12.8cm},\LenToUnit{6.3cm}){\waterword}%
  \put(\LenToUnit{18.1cm},\LenToUnit{4.9cm}){\waterword}%
  \put(\LenToUnit{-0.3cm},\LenToUnit{10.0cm}){\waterword}%
  \put(\LenToUnit{5.0cm},\LenToUnit{11.5cm}){\waterword}%
  \put(\LenToUnit{10.3cm},\LenToUnit{10.1cm}){\waterword}%
  \put(\LenToUnit{15.6cm},\LenToUnit{11.6cm}){\waterword}%
  \put(\LenToUnit{2.2cm},\LenToUnit{16.5cm}){\waterword}%
  \put(\LenToUnit{7.5cm},\LenToUnit{15.1cm}){\waterword}%
  \put(\LenToUnit{12.8cm},\LenToUnit{16.6cm}){\waterword}%
  \put(\LenToUnit{18.1cm},\LenToUnit{15.2cm}){\waterword}%
}
\usepackage{fancyhdr}
\usepackage{lastpage}
\usepackage{amsmath}
\usepackage{booktabs}
\usepackage{longtable}
\usepackage{xcolor}
\DeclareMathSizes{10.53937}{10.95}{7.5}{5.5}
\setlength{\parindent}{0pt}
\setlength{\parskip}{3pt}
\linespread{1.25}

\newcommand{\ans}[2]{\noindent\textbf{#1} \quad 答案：#2\par}
\newcommand{\expn}[1]{\hspace*{1.5em}解析：#1\par}
\newcommand{\bigtitle}[1]{\begin{center}\dengxian\large #1\end{center}}
\newcommand{\overview}[1]{\par\smallskip\noindent\textbf{程序总析：}#1\par\smallskip}
\newCJKfontfamily{\dengxian}{DengXian}

\pagestyle{fancy}
\fancyhf{}
\fancyfoot[C]{\xeCJKsetup{CJKecglue={\hskip0pt}}{\songti 折工 C5P-S 2026 参考答案与解析}\\[2pt]{\songti 第\thepage 页，共\pageref{LastPage}页}}
\renewcommand{\headrulewidth}{0pt}
\renewcommand{\footrulewidth}{0pt}

\begin{document}

\begin{center}
{\dengxian
{\fontsize{15.95}{1}\selectfont 2026 €€£超专业级别硬件能力认证第一轮}\\[1em]
{\fontsize{14.05}{1}\selectfont （C5P-S1）提高级 C++语言试题 \textbf{参考答案与解析}}\\[1em]
{\fontsize{12}{1}\selectfont 认证时间：2026年9月18日14:30\textasciitilde16:30}
}
\end{center}

\vspace{6pt}
\noindent\textbf{考生注意事项：}
\begin{itemize}
  \item 试题纸共有 15 页，答题纸共有 1 页，满分 100 分。请在答题纸上作答，写在试题纸上的一律无效。
  \item 不得使用任何电子设备（如计算器、手机、电子词典等）或查阅任何书籍资料。
\end{itemize}

\noindent\textbf{说明：}本文件对应自制模拟卷，不代表任何官方机构。题面、代码、解析和考点矩阵按本套 42 题版本同步。

\section*{分值与题号}
\begin{center}
\begin{tabular}{c c c c}
\toprule
部分 & 题号 & 题量与计分 & 小计\\
\midrule
单项选择 & 1--15 & $15\times2$ 分 & 30 分\\
阅读程序（1）KMP & 16--20 & $3\times1.5+1\times3+1\times4$ 分 & 11.5 分\\
阅读程序（2）状态压缩 DP & 21--26 & $2\times1.5+1\times2+3\times3$ 分 & 14 分\\
阅读程序（3）线段树 & 27--32 & $3\times1.5+2\times3+1\times4$ 分 & 14.5 分\\
完善程序 & 33--42 & $10\times3$ 分 & 30 分\\
\midrule
\multicolumn{3}{r}{总分} & 100 分\\
\bottomrule
\end{tabular}
\end{center}

\section*{一、单项选择题（1--15）}
\ans{1.}{B}\expn{\texttt{cat} 用于连接并显示文件内容。}
\ans{2.}{C}\expn{含 $m$ 个叶子的二叉树共有 $2m-1$ 个结点，指针域总数为 $4m-2$；非空指针为 $2m-2$，空指针域为 $2m$。}
\ans{3.}{B}\expn{完全二叉树高度为 $O(\log n)$，叶子数为 $O(n)$，逐个向根访问的总量为 $O(n\log n)$。}
\ans{4.}{D}\expn{这是长度为 10 的合法括号序列计数，答案为卡特兰数 $C_5=42$。}
\ans{5.}{A}\expn{只有在 B 为空时把 A 整体转移到 B，才能保证最早入队元素位于 B 栈顶。}
\ans{6.}{C}\expn{令 $y_i=x_i-(i-1)$，问题转化为从 7 个位置中选 4 个，答案为 $\binom74=35$。}
\ans{7.}{C}\expn{\texttt{set} 按键有序且不重复，基于平衡树时插入、删除、查找均为 $O(\log n)$。}
\ans{8.}{C}\expn{拓扑序要求每条有向边的起点排在终点之前；有环图不存在拓扑序。}
\ans{9.}{A}\expn{\texttt{next} 数组记录模式串失配后可以回退到的位置。}
\ans{10.}{A}\expn{朴素 Dijkstra 每轮线性寻找最小未确定顶点，总复杂度为 $O(n^2)$。}
\ans{11.}{D}\expn{模质数 $p$ 下，$1/k+1/(p-k)\equiv0$。括号内缺少 $1$，因此为 $-1$；指数 $p$ 为奇数，结果仍为 $p-1=20260912$。}
\ans{12.}{C}\expn{$1/(i(i+1))=1/i-1/(i+1)$，求和后为 $1-1/(p-1)\equiv2\pmod p$。}
\ans{13.}{B}\expn{枚举 9 条边中删除 2 条，保留至少一条从 1 到 8 的路径，可行方案共有 15 种。}
\ans{14.}{B}\expn{用容斥计算：$200-100-66-28+33+14+9-4=58$。}
\ans{15.}{C}\expn{由 $A(0,n)=n+1$、$A(1,n)=n+2$、$A(2,n)=2n+3$，可得 $A(2,3)=9$。}

\section*{二、阅读程序（16--32）}

\subsection*{（1）KMP 前缀函数}
\overview{程序先求每个前缀的最长真前后缀长度，再沿着前缀函数回退，累加整个字符串及各级相同前后缀的长度。}
\ans{16.}{A（正确）}\expn{输入为 \texttt{ababa} 时，各位置的前缀函数值为 \texttt{0,0,1,2,3}，所以 \texttt{pi[4]=3}。}
\ans{17.}{A（正确）}\expn{原循环第一次累加的是整个字符串长度 \texttt{n}，之后才沿 \texttt{pi} 回退。改为从 \texttt{pi[n-1]} 开始后只少累加这一项，因此输出恰好减少 \texttt{n}。}
\ans{18.}{B（错误）}\expn{一次失配后可能还需要继续沿着 \texttt{pi} 回退。把 \texttt{while} 改成 \texttt{if} 只回退一次，某些字符串的前缀函数会被算错，最终输出也可能改变。}
\ans{19.}{C}\expn{该字符串的前缀函数回退链给出的相同前后缀长度为 $14,8,2$，其和为 $24$。}
\ans{20.}{B}\expn{构造 \texttt{pi} 时，\texttt{j} 的增加和回退总次数都是线性级别；最后沿回退链访问的长度也不超过 $n$，所以时间复杂度为 $O(n)$。\texttt{pi} 数组和字符串占用 $O(n)$ 空间。}

\subsection*{（2）状态压缩 DP}
\overview{每一行用无相邻 1 的位掩码表示选取位置，转移要求相邻两行掩码不相交。}
\ans{21.}{A（正确）}\expn{三列中不出现相邻 1 的状态为 \texttt{000,001,010,100,101}，共 5 个。}
\ans{22.}{A（正确）}\expn{\texttt{101\&001=001}，结果非零，因此会执行 \texttt{continue}，不会进行这次转移。}
\ans{23.}{B（错误）}\expn{每轮写入的是当前行对应的滚动数组。如果不先清零，上一轮保留的数会被继续累加，某些输入的方案数会变大。}
\ans{24.}{B}\expn{预处理状态为 $O(2^m)$，每行枚举状态对，复杂度为 $O(nV^2+2^m)$。}
\ans{25.}{C}\expn{两行三列全可选时，按行状态转移统计得到 17。}
\ans{26.}{B}\expn{这两条语句把原本的两个下标交换回来，使 \texttt{pre} 指向上一轮、\texttt{cur} 指向本轮。删去后第一轮会清空原有的 \texttt{f[0][0]=1}，并从全为 0 的数组转移，之后每轮也都只能得到 0。}

\subsection*{（3）区间仿射更新线段树}
\overview{每个结点保存所代表区间的总和。一次完整覆盖的更新可以先记在当前结点，之后访问子区间时再传给两个子结点。}
\ans{27.}{A（正确）}\expn{\texttt{upd} 先判断当前区间是否被完整覆盖；若是，就调用一次 \texttt{apply} 修改当前结点，然后执行 \texttt{return}，不会继续递归。}
\ans{28.}{B（错误）}\expn{\texttt{add} 记录的加法部分也必须清零。否则同一个加法会在以后再次下传，可能被重复计算，导致查询结果错误。}
\ans{29.}{A（正确）}\expn{建树为 $O(n)$，每次更新或查询为 $O(\log n)$，总量可写成 $O(n+q\log n)$，也属于 $O((n+q)\log n)$。}
\ans{30.}{B}\expn{区间中共有 \texttt{r-l+1} 个数，每个数先乘 \texttt{k} 再加 \texttt{b}，所以总和变为 \texttt{tr[p].sum*k+(r-l+1)*b}，最后按模数取模。}
\ans{31.}{A}\expn{左右子结点互不相同，先更新左子结点还是先更新右子结点，都得到相同的结果；其余三项都会丢失或重复应用部分更新。}
\ans{32.}{B（4 分）}\expn{三次更新后的数组为 \texttt{8 11 27 39 13 13}，因此区间 $[2,6]$ 的和为 $11+27+39+13+13=103$。}

\section*{三、完善程序（33--42）}
\subsection*{（1）三色序列}
\ans{33.}{C}\expn{先把前一个位置的计数复制过来；如果当前字符属于颜色 \texttt{t}，再多计数一次，因此填 \texttt{pre[t][i-1] + (t == c)}。}
\ans{34.}{B}\expn{\texttt{f} 记录的是交换次数，初始状态尚未构造任何字符，代价应为 0。}
\ans{35.}{A}\expn{前 \texttt{len} 个位置已经构造，其中使用了 \texttt{r} 个 R 和 \texttt{g} 个 G，所以 Y 的使用数为 \texttt{len-r-g}。}
\ans{36.}{A}\expn{新字符不能和上一字符相同，也不能超过该颜色的总数；满足任一条件就跳过，因此填 \texttt{c == last || used[c] == cnt[c]}。}
\ans{37.}{C}\expn{当前状态代价为 \texttt{val}，加入该字符的增量为 \texttt{w}，所以新值为 \texttt{val + w}。}

\subsection*{（2）盒子}
\ans{38.}{B}\expn{\texttt{lo} 保存当前选中的较小间隙；若数量过多，应移出其中最大的一个，因此填 \texttt{prev(lo.end())}。}
\ans{39.}{B}\expn{\texttt{x} 是 \texttt{lo} 中最大的间隙，\texttt{y} 是 \texttt{hi} 中最小的间隙。只有在 \texttt{*x <= *y} 时才不需要交换。}
\ans{40.}{B}\expn{交换后，\texttt{lo} 中的 \texttt{vx} 被 \texttt{vy} 替换，答案变化为 \texttt{vy-vx}；代码用 \texttt{ans +=} 累加这个变化量。}
\ans{41.}{A}\expn{\texttt{n} 个盒子最终形成 \texttt{k} 个链，每条链需要相邻盒子之间的间隙，共需保留 \texttt{n-k} 个间隙。}
\ans{42.}{C}\expn{删除体积为 \texttt{v} 的中间盒子后，左、右盒子直接相邻，新间隙为 \texttt{r-l}。}

\end{document}
